\section{Cơ sở lý thuyết}

\subsection{Lý thuyết công nghệ sử dụng}

\subsubsection{Mô hình Ngôn ngữ Lớn (Large Language Model)}

\textbf{Định nghĩa:} Mô hình Ngôn ngữ Lớn (LLM) là các mạng nơ-ron sâu với
hàng tỷ tham số, được huấn luyện trên tập văn bản khổng lồ để học phân phối
xác suất của ngôn ngữ tự nhiên. Kiến trúc nền tảng là \textbf{Transformer}
\cite{vaswani2017attention} với cơ chế \textit{self-attention} cho phép nắm
bắt phụ thuộc xa trong chuỗi văn bản.

\textbf{Cơ chế hoạt động:} Khi nhận một đầu vào (prompt), LLM sinh ra đầu ra
bằng cách dự đoán token tiếp theo lặp đi lặp lại (autoregressive decoding).
Tham số \texttt{temperature} điều chỉnh tính ngẫu nhiên: giá trị thấp ($<0.3$)
cho đầu ra nhất quán, giá trị cao ($>0.7$) cho đầu ra sáng tạo hơn.

\textbf{Ứng dụng trong đề tài:} LLM được sử dụng cho hai tác vụ:
\begin{itemize}
    \item \textbf{Tóm tắt (Summarization):} Từ một chunk văn bản, sinh ra
          danh sách 3--5 ý chính ngắn gọn.
    \item \textbf{Sinh câu hỏi (Question Generation):} Từ cùng chunk, sinh
          ra 5--10 câu hỏi trắc nghiệm 4 đáp án kèm giải thích.
\end{itemize}

\subsubsection{Prompt Engineering}

\textbf{Định nghĩa:} Prompt Engineering là kỹ thuật thiết kế đầu vào cho LLM
để dẫn hướng đầu ra về chất lượng, định dạng và nội dung mong muốn, mà không
cần thay đổi tham số mô hình (zero-shot / few-shot).

\textbf{Các kỹ thuật áp dụng trong đề tài:}

\begin{table}[H]
    \centering
    \caption{Các kỹ thuật Prompt Engineering và ứng dụng trong đề tài.}
    \label{tab:prompt-tech}
    \begin{tabular}{p{3.5cm}p{6cm}p{5cm}}
        \toprule
        \textbf{Kỹ thuật} & \textbf{Mô tả} & \textbf{Ứng dụng} \\
        \midrule
        \textbf{Role Prompting} &
            Gán vai trò chuyên gia cho LLM để định hướng phong cách. &
            ``Bạn là giảng viên đại học soạn đề kiểm tra\ldots'' \\
        \midrule
        \textbf{Structured Output} &
            Yêu cầu LLM trả về JSON theo schema cố định. &
            \texttt{response\_format = json\_object}; schema kiểm tra bằng Pydantic. \\
        \midrule
        \textbf{Constraint Prompting} &
            Đặt ràng buộc rõ ràng về số lượng, độ dài, ngôn ngữ. &
            ``Tạo đúng \{n\} câu hỏi; mỗi câu $\leq$ 30 từ; ngôn ngữ: tiếng Việt.'' \\
        \midrule
        \textbf{Few-shot Examples} &
            Cung cấp ví dụ mẫu để định hướng format. &
            Đính kèm 1 ví dụ câu hỏi mẫu trong prompt. \\
        \bottomrule
    \end{tabular}
\end{table}

\subsubsection{Text Chunking}

\textbf{Định nghĩa:} Text Chunking là quá trình chia văn bản dài thành các
đoạn nhỏ hơn (chunk) có tính ngữ nghĩa độc lập, phục vụ cho xử lý LLM và
lưu trữ vector.

\textbf{Vì sao cần chunking?} LLM có giới hạn \textit{context window}
(ví dụ GPT-4o-mini: 128K token). Tuy nhiên, đưa toàn bộ tài liệu vào một
lần gọi API vừa tốn kém vừa làm giảm chất lượng đầu ra do LLM khó tập trung
vào từng phần. Chunking giải quyết vấn đề này bằng cách xử lý từng đoạn nhỏ.

\textbf{Bài toán thiết kế:} Chunk lý tưởng cần thỏa mãn ba điều kiện:
(1)~tính ngữ nghĩa (một ý hoàn chỉnh), (2)~kích thước phù hợp (500--1500
token), (3)~metadata đầy đủ (thuộc section nào, trang bao nhiêu).

\subsubsection{Vector Embedding}

\textbf{Định nghĩa:} Vector Embedding là ánh xạ $f: \text{text} \to
\mathbb{R}^d$ sao cho các văn bản có nghĩa tương đồng có véc-tơ gần nhau
theo độ đo cosine:
\[
    \text{sim}(\vec{u}, \vec{v}) = \frac{\vec{u} \cdot \vec{v}}{|\vec{u}||\vec{v}|}
\]

\textbf{Ứng dụng:} Mỗi chunk được chuyển thành véc-tơ $\vec{e} \in
\mathbb{R}^{1024}$ và lưu vào Qdrant. Khi học viên tìm kiếm, query cũng được
chuyển thành véc-tơ và Qdrant tìm top-K chunk có similarity cao nhất --- đây
là nền tảng của tìm kiếm ngữ nghĩa và RAG.

\subsubsection{Domain-Driven Design (DDD)}

\textbf{Định nghĩa:} DDD \cite{evans2004ddd} là phương pháp thiết kế phần
mềm tổ chức code xung quanh mô hình nghiệp vụ. Nguyên tắc cốt lõi:
logic nghiệp vụ phải độc lập với framework và hạ tầng kỹ thuật.

\textbf{Áp dụng trong đề tài --- bốn lớp (Hexagonal Architecture):}
\begin{enumerate}
    \item \textbf{Domain:} Entity (\texttt{Document}, \texttt{Chunk},
          \texttt{LearningOutcome}, \texttt{Curriculum}), Port interface khai
          báo contract (\texttt{ILLMClient}, \texttt{IDocumentParser},
          \texttt{IMetadataStore}, \texttt{IGraphStore}, \texttt{ICurriculumExtractor}).
    \item \textbf{Application:} Use case (\texttt{ProcessDocumentUseCase},
          \texttt{IngestCurriculumUseCase}, \texttt{SearchByLearningOutcomeUseCase},
          \texttt{GenerateCurriculumQuizUseCase}) --- chỉ phụ thuộc Port.
    \item \textbf{Infrastructure:} Adapter cài đặt Port (\texttt{GeminiAdapter},
          \texttt{PdfParser}, \texttt{PostgresMetadataStore},
          \texttt{Neo4jGraphStore}, \texttt{QdrantVectorStore}, \texttt{MinIOAdapter}).
    \item \textbf{Delivery:} gRPC Servicer nhận request, gọi use case; Worker
          tiêu thụ job từ Redis queue.
\end{enumerate}

\subsubsection{Knowledge Graph và Chuẩn đầu ra học phần (LO)}

\textbf{Định nghĩa:} Knowledge Graph (KG) là cơ sở tri thức biểu diễn các
thực thể và quan hệ giữa chúng dưới dạng đồ thị có hướng và có nhãn.
Trong đề tài, KG được dùng để biểu diễn cấu trúc môn học theo chuẩn CDIO
của ĐHBK TP.HCM.

\textbf{Cấu trúc KG môn học:} Mỗi môn học bao gồm:
\begin{itemize}
    \item Nút \textbf{Course} --- thông tin môn học (mã môn, tên, tín chỉ).
    \item Nút \textbf{Chapter} --- chương bài giảng (có thứ tự).
    \item Nút \textbf{LearningOutcome (LO)} --- Chuẩn đầu ra học phần,
          phân cấp cha--con (ví dụ: L.O.3 $\to$ L.O.3.1).
    \item Nút \textbf{Assessment} --- Hình thức đánh giá (giữa kỳ, cuối kỳ).
    \item Cạnh \textbf{SUPPORTS} --- từ Chunk nội dung đến LO, kèm trọng số
          (confidence).
\end{itemize}

\textbf{Ứng dụng:} Khi giảng viên yêu cầu tạo Quiz cho một LO hoặc một
chương, hệ thống duyệt KG để lấy các chunk nội dung liên quan (thông qua cạnh
SUPPORTS) và đưa chúng vào context cho LLM --- đảm bảo Quiz được gắn với
đúng Chuẩn đầu ra thay vì chỉ dựa trên similarity vector thuần túy.

\subsubsection{GraphRAG: Truy xuất Tăng cường dựa trên Đồ thị}

\textbf{Định nghĩa:} GraphRAG (Graph-based Retrieval-Augmented Generation) là
phương pháp kết hợp Knowledge Graph với RAG truyền thống nhằm khắc phục hạn
chế của việc tìm kiếm thuần vector-similarity. Thay vì chỉ so khớp đoạn văn
gần nhất về mặt ngữ nghĩa, GraphRAG duyệt các quan hệ có cấu trúc giữa các
thực thể trong đồ thị tri thức để truy xuất ngữ cảnh chính xác hơn.

\textbf{Hạn chế của RAG truyền thống:}
\begin{itemize}
    \item \textbf{Chunking cơ học:} Cắt văn bản theo kích thước cố định làm
          đứt gãy ngữ cảnh giữa các đoạn liền kề.
    \item \textbf{Thiếu tư duy tổng hợp:} Câu hỏi ở mức độ \textit{Vận dụng}
          và \textit{Phân tích} (Bloom) đòi hỏi liên kết nhiều khái niệm ---
          RAG thuần túy không mô hình hóa được các liên kết này.
    \item \textbf{Ảo giác (Hallucination):} Khi không tìm được đoạn phù hợp,
          LLM có xu hướng bịa thông tin.
\end{itemize}

\textbf{Giải pháp GraphRAG:} Văn bản thô được chuyển đổi thành các bộ ba cấu
trúc (Triplets) $\langle$Chủ thể, Mối quan hệ, Khách thể$\rangle$. Khi cần
sinh câu hỏi, hệ thống duyệt đồ thị từ nút khái niệm đích, thu thập các nút
lân cận liên quan, và đưa ngữ cảnh có cấu trúc này vào LLM --- thay vì chỉ
cung cấp các đoạn văn tương đồng theo cosine.

\begin{table}[H]
    \centering
    \caption{So sánh RAG truyền thống và GraphRAG.}
    \label{tab:rag-vs-graphrag}
    \begin{tabular}{p{4cm}p{5cm}p{5cm}}
        \toprule
        \textbf{Tiêu chí} & \textbf{RAG truyền thống} & \textbf{GraphRAG} \\
        \midrule
        Đơn vị lưu trữ     & Chunk văn bản        & Node + Edge (Triplet) \\
        Tìm kiếm            & Cosine similarity    & Graph traversal + similarity \\
        Liên kết khái niệm  & Ngầm định            & Tường minh (có nhãn) \\
        Câu hỏi tổng hợp    & Khó                  & Tự nhiên \\
        Chi phí xây dựng    & Thấp                 & Cao hơn (cần Graph Extraction) \\
        \bottomrule
    \end{tabular}
\end{table}

\noindent\textbf{Ứng dụng trong đề tài:} GraphRAG được áp dụng để sinh câu
hỏi ở mức độ cao theo thang Bloom (Vận dụng, Phân tích). Các node trong đồ
thị ánh xạ tới các khái niệm chủ chốt; các edge mã hóa mối quan hệ ngữ nghĩa
(ví dụ: \textit{``Overfitting''} $\xrightarrow{\text{giải pháp}}$
\textit{``Regularization''}). LLM nhận subgraph thay vì đoạn văn thuần túy,
cho phép đặt câu hỏi đòi hỏi lập luận chéo.

\subsection{Cân nhắc lựa chọn công nghệ}

\subsubsection{Lựa chọn LLM API: GPT-4o-mini vs Gemini Flash}

\begin{table}[H]
    \centering
    \caption{So sánh GPT-4o-mini và Gemini Flash.}
    \label{tab:llm-choice}
    \begin{tabular}{p{4cm}p{5cm}p{5.5cm}}
        \toprule
        \textbf{Tiêu chí} & \textbf{GPT-4o-mini} & \textbf{Gemini 2.0 Flash} \\
        \midrule
        Chất lượng tiếng Việt & Tốt & Tốt (đa ngôn ngữ tự nhiên) \\
        Chi phí (input/1M token) & \$0.15 & \$0.10 (free tier rộng) \\
        Tốc độ                & Nhanh ($\sim$30 token/s) & Rất nhanh ($\sim$60 token/s) \\
        JSON mode             & Có (\texttt{json\_object}) & Có (\texttt{application/json}) \\
        Độ ổn định API        & Rất cao & Cao \\
        \bottomrule
    \end{tabular}
\end{table}

\noindent\textbf{Quyết định:} Sử dụng \textbf{Gemini 2.0 Flash} làm mặc định do
chi phí thấp hơn, tốc độ nhanh hơn, và free tier phù hợp cho môi trường nghiên
cứu. Kiến trúc Port/Adapter cho phép chuyển sang GPT-4o-mini dễ dàng nếu cần.

\subsubsection{Chiến lược Hybrid LLMs: Open-weights kết hợp Closed API}

Để cân bằng giữa chi phí, bảo mật và hiệu năng, hệ thống áp dụng chiến lược
kết hợp hai loại mô hình ngôn ngữ với vai trò phân công rõ ràng:

\begin{table}[H]
    \centering
    \caption{Phân công vai trò trong chiến lược Hybrid LLMs.}
    \label{tab:hybrid-llm}
    \begin{tabular}{p{3.5cm}p{5.5cm}p{5.5cm}}
        \toprule
        \textbf{Tiêu chí} & \textbf{Open-weights (Local)} & \textbf{Closed API (Cloud)} \\
        \midrule
        Ví dụ mô hình      & Gemma 3, Llama 3.3            & Gemini 2.5 Flash, GPT-4o \\
        Triển khai         & Cục bộ / On-premise           & Qua API \\
        Bảo mật dữ liệu    & Tuyệt đối (offline)           & Phụ thuộc chính sách provider \\
        Tinh chỉnh         & Có (LoRA/QLoRA)               & Không \\
        Vai trò trong pipeline & Graph Extraction (trích xuất triplet) & Sinh câu hỏi \& Đánh giá chất lượng \\
        \bottomrule
    \end{tabular}
\end{table}

\noindent\textbf{Lý do phân công:} Tác vụ Graph Extraction chạy trên toàn bộ
nội dung tài liệu (nhiều token, lặp lại nhiều lần) nên ưu tiên mô hình local
để kiểm soát chi phí và bảo mật. Tác vụ sinh câu hỏi yêu cầu lập luận mạnh,
tốc độ cao và chất lượng đầu ra tốt --- phù hợp với các Closed API hiện đại.

\subsubsection{Lựa chọn chiến lược Chunking}

\begin{table}[H]
    \centering
    \caption{So sánh các chiến lược Chunking cho tài liệu học thuật.}
    \label{tab:chunking-choice}
    \begin{tabular}{p{3cm}p{5.5cm}p{6cm}}
        \toprule
        \textbf{Chiến lược} & \textbf{Ưu điểm} & \textbf{Nhược điểm} \\
        \midrule
        Fixed-size & Đơn giản; kích thước nhất quán. & Cắt đứt câu/ý; mất ngữ cảnh. \\
        Sentence-based & Bảo toàn câu hoàn chỉnh. & Chunk nhỏ; thiếu ngữ cảnh đoạn. \\
        \textbf{Heading-based} & Phù hợp tài liệu có cấu trúc; mỗi chunk = 1 chủ đề. & Phụ thuộc chất lượng cấu trúc tài liệu gốc. \\
        Semantic & Chính xác nhất. & Chi phí tính toán cao. \\
        \bottomrule
    \end{tabular}
\end{table}

\noindent\textbf{Quyết định:} \textbf{Heading-based} là chiến lược chính.
Section vượt ngưỡng sẽ được chia tiếp bằng sliding window với overlap 100
token để bảo toàn ngữ cảnh.

\subsubsection{Lựa chọn Vector Database: Qdrant vs Weaviate vs Pinecone}

\begin{table}[H]
    \centering
    \caption{So sánh các Vector Database.}
    \label{tab:vdb-choice}
    \begin{tabular}{p{3cm}p{3.5cm}p{3.5cm}p{4.5cm}}
        \toprule
        \textbf{Tiêu chí} & \textbf{Qdrant} & \textbf{Weaviate} & \textbf{Pinecone} \\
        \midrule
        Tự host (on-premise)  & Có & Có & Không (cloud-only) \\
        Hiệu năng ANN         & Rất cao (HNSW) & Cao & Cao \\
        Payload filtering     & Có & Có & Có (giới hạn) \\
        Python SDK            & Đơn giản & Phức tạp hơn & Đơn giản \\
        Chi phí               & Miễn phí (OSS) & Miễn phí (OSS) & Có phí (theo vector) \\
        \bottomrule
    \end{tabular}
\end{table}

\noindent\textbf{Quyết định:} \textbf{Qdrant} --- mã nguồn mở, tự host,
hiệu năng cao, Python SDK thân thiện.

\subsubsection{Lựa chọn Embedding Model: bge-m3 vs text-embedding-3-small}

\begin{table}[H]
    \centering
    \caption{So sánh mô hình Embedding.}
    \label{tab:embed-choice}
    \begin{tabular}{p{4cm}p{5cm}p{5.5cm}}
        \toprule
        \textbf{Tiêu chí} & \textbf{BAAI/bge-m3} & \textbf{text-embedding-3-small} \\
        \midrule
        Chiều véc-tơ     & 1024 & 1536 \\
        Hỗ trợ tiếng Việt & Rất tốt (đa ngôn ngữ) & Tốt \\
        Chi phí           & Miễn phí (local) & \$0.02/1M token \\
        Tốc độ inference  & Nhanh (CPU đủ dùng) & Rất nhanh (API call) \\
        \bottomrule
    \end{tabular}
\end{table}

\noindent\textbf{Quyết định:} \textbf{BAAI/bge-m3} chạy local --- không tốn
chi phí API, hỗ trợ tiếng Việt tốt, phù hợp môi trường staging.

\subsubsection{Lựa chọn Graph Database: Neo4j}

\textbf{Lý do chọn Neo4j:} Cấu trúc Course $\to$ Chapter $\to$ LO $\to$ Chunk
là đồ thị phân cấp có nhãn, phù hợp tự nhiên với mô hình đồ thị. Truy vấn
``Tìm tất cả chunk hỗ trợ LO.3.1'' viết bằng Cypher ngắn gọn hơn JOIN nhiều
bảng trong SQL. Neo4j là lựa chọn mã nguồn mở, tự host, với Python driver
\texttt{neo4j} chính thức.

\subsubsection{Lựa chọn nền tảng LMS tích hợp: Open edX vs Moodle}

Hệ thống cần tích hợp với một nền tảng LMS hiện có để quản lý người dùng,
khóa học và tiến độ học tập. Hai lựa chọn được đánh giá là \textbf{Open edX}
(triển khai qua Tutor) và \textbf{Moodle}.

\paragraph{Hạ tầng \& Vận hành}

\begin{table}[H]
    \centering
    \caption{So sánh hạ tầng và vận hành giữa Open edX và Moodle.}
    \label{tab:edx-moodle-infra}
    \begin{tabular}{p{4.5cm}p{4.5cm}p{5cm}}
        \toprule
        \textbf{Tiêu chí} & \textbf{Open edX (Tutor)} & \textbf{Moodle} \\
        \midrule
        RAM tối thiểu           & $\geq$ 4--6 GB                           & $\geq$ 1--2 GB \\
        Thời gian cài đặt       & $\sim$30--60 phút                        & $\sim$15--30 phút \\
        Cài Plugin/XBlock       & CLI (\texttt{tutor plugins} + \texttt{pip}) & Upload \texttt{.zip} qua Admin UI \\
        CI/CD                   & Phức tạp (build custom image)            & Đơn giản \\
        Cập nhật phiên bản      & Theo release (Sumac, Redwood\ldots)      & Tự động / dễ dàng \\
        \bottomrule
    \end{tabular}
\end{table}

\paragraph{Plugin \& Khả năng mở rộng}

\begin{table}[H]
    \centering
    \caption{So sánh khả năng mở rộng giữa Open edX và Moodle.}
    \label{tab:edx-moodle-plugin}
    \begin{tabular}{p{4.5cm}p{4.5cm}p{5cm}}
        \toprule
        \textbf{Tiêu chí} & \textbf{Open edX} & \textbf{Moodle} \\
        \midrule
        Cơ chế mở rộng             & XBlock (Python)           & Plugin (PHP, AMD JS + Mustache) \\
        Chuẩn liên thông LTI       & LTI 1.3 (chỉ nhận)        & LTI 1.3 (nhận \textbf{và} phát) \\
        Tích hợp app ngoài no-code & Hạn chế (cần code Python) & Linh hoạt (iframe qua LTI) \\
        Độ khó phát triển          & Trung bình (cần hiểu XBlock) & Cao hơn nếu yếu PHP \\
        \bottomrule
    \end{tabular}
\end{table}

\noindent\textbf{Nhận xét:} Moodle có ưu thế về tài nguyên thấp, cài đặt đơn
giản và hỗ trợ LTI 1.3 hai chiều --- phù hợp với chiến lược tích hợp module
AI dưới dạng \textit{LTI Tool Provider} độc lập. Open edX phù hợp hơn khi đội
ngũ có nền tảng Python mạnh và hạ tầng server lớn. Quyết định cuối cùng sẽ
được xác nhận sau khi tham khảo ý kiến GVHD về môi trường triển khai thực tế.

\subsection{Tổng kết Chương 3}

Chương 3 đã trình bày nền tảng lý thuyết cho các công nghệ cốt lõi của hệ
thống: LLM và Prompt Engineering, Text Chunking, Vector Embedding, kiến trúc
DDD (Hexagonal Architecture), Knowledge Graph cho Chuẩn đầu ra học phần, và
GraphRAG --- phương pháp truy xuất có cấu trúc vượt trội so với RAG thuần
vector. Phần cân nhắc lựa chọn công nghệ đã lý giải cụ thể lý do chọn chiến
lược Hybrid LLMs, Gemini 2.0 Flash, Heading-based Chunking, Qdrant, bge-m3,
Neo4j, và so sánh hai nền tảng LMS tích hợp (Open edX vs Moodle).
Những quyết định này được hiện thực hóa trong Chương 4 và 5.
